% =====================================================================
%  Annexes (Buttler) : éléments externes au rapport, lecture optionnelle.
% =====================================================================

\chapter{Liste des hyperparamètres optimisés}
\label{ann:hyperparams}

Le tableau suivant détaille l'espace de recherche utilisé par Optuna en phase~3.

\begin{table}[H]
  \centering
  \caption{Espace de recherche Optuna.}
  \label{tab:ann:optuna}
  \begin{tabular}{lll}
    \toprule
    \textbf{Hyperparamètre} & \textbf{Plage} & \textbf{Type} \\
    \midrule
    Taille de fenêtre \(W\)        & [30, 60, 90, 120] & catégoriel \\
    Filtres conv. couche 1         & [16, 32, 64]      & catégoriel \\
    Filtres conv. couche 2         & [32, 64, 128]     & catégoriel \\
    Unités LSTM                    & [32, 64, 128]     & catégoriel \\
    Dropout                        & [0,1 ; 0,5]       & continu (log) \\
    Taux d'apprentissage           & [1e-5 ; 1e-3]     & continu (log) \\
    Taille de \textit{batch}       & [32, 64, 128]     & catégoriel \\
    \bottomrule
  \end{tabular}
\end{table}

\chapter{Reproductibilité et code source}
\label{ann:reproductibilite}

Le code complet du projet est disponible publiquement à l'adresse suivante :

\medskip
\centerline{\url{[À COMPLÉTER : URL du dépôt GitHub]}}
\medskip

Le fichier \texttt{requirements-lock.txt} fixe les versions exactes des bibliothèques utilisées. La commande \texttt{make demo} lance l'application interactive en local après installation des dépendances.

\chapter{Captures d'écran complémentaires}
\label{ann:screenshots}

\begin{figure}[H]
  \centering
  \figplaceholder[0.8\linewidth][5cm]{Capture — onglet Model Showdown}
  \caption{Onglet \emph{Model Showdown} de l'application interactive.}
\end{figure}

\begin{figure}[H]
  \centering
  \figplaceholder[0.8\linewidth][5cm]{Capture — onglet Compression Lab}
  \caption{Onglet \emph{Compression Lab}.}
\end{figure}

\begin{figure}[H]
  \centering
  \figplaceholder[0.8\linewidth][5cm]{Capture — onglet What-If}
  \caption{Onglet \emph{What-If}.}
\end{figure}
